\section{More on images}
\label{section-more-images}

\noindent
In this section we collect a few additional lemmas concerning the image
on $\Spec$ for ring maps. See also Section \ref{section-going-up}
for example.

\begin{lemma}
\label{lemma-generic-finite-presentation}
Let $R \subset S$ be an inclusion of domains.
Assume that $R \to S$ is of finite type.
There exists a nonzero $f \in R$, and a nonzero $g \in S$
such that $R_f \to S_{fg}$ is of finite presentation.
\end{lemma}

\begin{proof}
By induction on the number of generators of $S$ over $R$.
During the proof we may replace $R$ by $R_f$ and $S$ by $S_f$
for some nonzero $f \in R$.

\medskip\noindent
Suppose that $S$ is generated by a single element
over $R$. Then $S = R[x]/\mathfrak q$ for some
prime ideal $\mathfrak q \subset R[x]$. If $\mathfrak q = (0)$
there is nothing to prove. If $\mathfrak q \not = (0)$,
then let $h \in \mathfrak q$ be a nonzero element with minimal
degree in $x$. Write $h = f x^d + a_{d - 1} x^{d - 1} + \ldots + a_0$
with $a_i \in R$ and $f \not = 0$. After inverting $f$
in $R$ and $S$ we may assume that $h$ is monic. We obtain
a surjective $R$-algebra map $R[x]/(h) \to S$.
We have $R[x]/(h) = R \oplus Rx \oplus \ldots \oplus Rx^{d - 1}$
as an $R$-module and by minimality of $d$ we see that
$R[x]/(h)$ maps injectively into $S$. Thus $R[x]/(h) \cong S$
is finitely presented over $R$.

\medskip\noindent
Suppose that $S$ is generated by $n > 1$ elements over $R$.
Say $x_1, \ldots, x_n \in S$ generate $S$. Denote $S' \subset S$
the $R$-subalgebra generated by $x_1, \ldots, x_{n-1}$. By induction
hypothesis we see that there exist $f\in R$ and $g \in S'$
nonzero such that $R_f \to S'_{fg}$ is of finite presentation.
Next we apply the induction hypothesis to $S'_{fg} \to S_{fg}$
to see that there exist $f' \in S'_{fg}$ and
$g' \in S_{fg}$ such that $S'_{fgf'} \to S_{fgf'g'}$
is of finite presentation. We leave it to the reader to conclude.
\end{proof}

\begin{lemma}
\label{lemma-characterize-image-finite-type}
Let $R \to S$ be a finite type ring map.
Denote $X = \Spec(R)$ and $Y = \Spec(S)$.
Write $f : Y \to X$ the induced
map of spectra. Let $E \subset Y = \Spec(S)$ be a
constructible set.
If a point $\xi \in X$ is in $f(E)$, then
$\overline{\{\xi\}} \cap f(E)$ contains an open
dense subset of $\overline{\{\xi\}}$.
\end{lemma}

\begin{proof}
Let $\xi \in X$ be a point of $f(E)$. Choose a point $\eta \in E$
mapping to $\xi$. Let $\mathfrak p \subset R$ be the prime
corresponding to $\xi$ and let $\mathfrak q \subset S$ be the
prime corresponding to $\eta$. Consider the diagram
$$
\xymatrix{
\eta \ar[r] \ar@{|->}[d] & E \cap Y' \ar[r] \ar[d] &
Y' = \Spec(S/\mathfrak q) \ar[r] \ar[d] &
Y \ar[d] \\
\xi \ar[r] & f(E) \cap X' \ar[r] &
X' = \Spec(R/\mathfrak p) \ar[r] &
X
}
$$
By Lemma \ref{lemma-affine-map-quasi-compact} the set $E \cap Y'$
is constructible in $Y'$.
It follows that we may replace $X$ by $X'$ and
$Y$ by $Y'$. Hence we may assume that $R \subset S$ is an
inclusion of domains, $\xi$ is the generic
point of $X$, and $\eta$ is the generic point of $Y$.
By Lemma \ref{lemma-generic-finite-presentation}
combined with Chevalley's theorem
(Theorem \ref{theorem-chevalley})
we see that there exist dense opens $U \subset X$,
$V \subset Y$ such that $f(V) \subset U$ and
such that $f : V \to U$ maps constructible sets
to constructible sets. Note that $E \cap V$ is
constructible in $V$, see Topology,
Lemma \ref{topology-lemma-open-immersion-constructible-inverse-image}.
Hence $f(E \cap V)$ is constructible in $U$ and contains $\xi$.
By Topology, Lemma \ref{topology-lemma-generic-point-in-constructible}
we see that $f(E \cap V)$ contains a dense open $U' \subset U$.
\end{proof}

\noindent
At the end of this section we present a few more results on
images of maps on Spectra that have nothing to do with constructible
sets.

\begin{lemma}
\label{lemma-surjective-spec-radical-ideal}
Let $\varphi : R \to S$ be a ring map.
The following are equivalent:
\begin{enumerate}
\item The map $\Spec(S) \to \Spec(R)$ is surjective.
\item For any ideal $I \subset R$
the inverse image of $\sqrt{IS}$ in $R$ is equal to $\sqrt{I}$.
\item For any radical ideal $I \subset R$ the inverse image
of $IS$ in $R$ is equal to $I$.
\item For every prime $\mathfrak p$ of $R$ the inverse
image of $\mathfrak p S$ in $R$ is $\mathfrak p$.
\end{enumerate}
In this case the same is true after any base change: Given a ring map
$R \to R'$ the ring map $R' \to R' \otimes_R S$ has the equivalent
properties (1), (2), (3) as well.
\end{lemma}

\begin{proof}
If $J \subset S$ is an ideal, then
$\sqrt{\varphi^{-1}(J)} = \varphi^{-1}(\sqrt{J})$. This shows that (2)
and (3) are equivalent.
The implication (3) $\Rightarrow$ (4) is immediate.
If $I \subset R$ is a radical ideal, then
Lemma \ref{lemma-Zariski-topology}
guarantees that $I = \bigcap_{I \subset \mathfrak p} \mathfrak p$.
Hence (4) $\Rightarrow$ (2). By
Lemma \ref{lemma-in-image}
we have $\mathfrak p = \varphi^{-1}(\mathfrak p S)$ if and only if
$\mathfrak p$ is in the image. Hence (1) $\Leftrightarrow$ (4).
Thus (1), (2), (3), and (4) are equivalent.

\medskip\noindent
Assume (1) holds. Let $R \to R'$ be a ring map. Let
$\mathfrak p' \subset R'$ be a prime ideal lying over the prime
$\mathfrak p$ of $R$. To see that $\mathfrak p'$ is in the image
of $\Spec(R' \otimes_R S) \to \Spec(R')$ we have to show
that $(R' \otimes_R S) \otimes_{R'} \kappa(\mathfrak p')$ is not zero, see
Lemma \ref{lemma-in-image}.
But we have
$$
(R' \otimes_R S) \otimes_{R'} \kappa(\mathfrak p') =
S \otimes_R \kappa(\mathfrak p)
\otimes_{\kappa(\mathfrak p)} \kappa(\mathfrak p')
$$
which is not zero as $S \otimes_R \kappa(\mathfrak p)$ is not zero
by assumption and $\kappa(\mathfrak p) \to \kappa(\mathfrak p')$ is
an extension of fields.
\end{proof}

\begin{lemma}
\label{lemma-domain-image-dense-set-points-generic-point}
Let $R$ be a domain. Let $\varphi : R \to S$ be a ring map.
The following are equivalent:
\begin{enumerate}
\item The ring map $R \to S$ is injective.
\item The image $\Spec(S) \to \Spec(R)$
contains a dense set of points.
\item There exists a prime ideal $\mathfrak q \subset S$
whose inverse image in $R$ is $(0)$.
\end{enumerate}
\end{lemma}

\begin{proof}
Let $K$ be the field of fractions of the domain $R$.
Assume that $R \to S$ is injective. Since localization
is exact we see that $K \to S \otimes_R K$ is injective.
Hence there is a prime mapping to $(0)$ by
Lemma \ref{lemma-in-image}.

\medskip\noindent
Note that $(0)$ is dense in $\Spec(R)$, so that the
last condition implies the second.

\medskip\noindent
Suppose the second condition holds. Let $f \in R$,
$f \not = 0$. As $R$ is a domain we see that $V(f)$
is a proper closed subset of $R$. By assumption
there exists a prime $\mathfrak q$
of $S$ such that $\varphi(f) \not \in \mathfrak q$.
Hence $\varphi(f) \not = 0$.
Hence $R \to S$ is injective.
\end{proof}

\begin{lemma}
\label{lemma-injective-minimal-primes-in-image}
Let $R \subset S$ be an injective ring map.
Then $\Spec(S) \to \Spec(R)$
hits all the minimal primes.
\end{lemma}

\begin{proof}
Let $\mathfrak p \subset R$ be a minimal prime.
In this case $R_{\mathfrak p}$ has a unique prime ideal.
Hence it suffices to show that $S_{\mathfrak p}$ is not zero.
And this follows from the fact that localization is exact,
see Proposition \ref{proposition-localization-exact}.
\end{proof}

\begin{lemma}
\label{lemma-image-dense-generic-points}
Let $R \to S$ be a ring map. The following are equivalent:
\begin{enumerate}
\item The kernel of $R \to S$ consists of nilpotent elements.
\item The minimal primes of $R$ are in the image of
$\Spec(S) \to \Spec(R)$.
\item The image of $\Spec(S) \to \Spec(R)$ is dense
in $\Spec(R)$.
\end{enumerate}
\end{lemma}

\begin{proof}
Let $I = \Ker(R \to S)$. Note that
$\sqrt{(0)} = \bigcap_{\mathfrak q \subset S} \mathfrak q$, see
Lemma \ref{lemma-Zariski-topology}.
Hence $\sqrt{I} = \bigcap_{\mathfrak q \subset S} R \cap \mathfrak q$.
Thus $V(I) = V(\sqrt{I})$ is the closure of the image of
$\Spec(S) \to \Spec(R)$.
This shows that (1) is equivalent to (3). It is clear that
(2) implies (3). Finally, assume (1). We may replace
$R$ by $R/I$ and $S$ by $S/IS$ without affecting the topology
of the spectra and the map. Hence the implication (1) $\Rightarrow$ (2)
follows from Lemma \ref{lemma-injective-minimal-primes-in-image}.
\end{proof}

\begin{lemma}
\label{lemma-minimal-prime-image-minimal-prime}
Let $R \to S$ be a ring map. If a minimal prime $\mathfrak p \subset R$
is in the image of $\Spec(S) \to \Spec(R)$, then it is the image
of a minimal prime.
\end{lemma}

\begin{proof}
Say $\mathfrak p = \mathfrak q \cap R$. Then choose a minimal
prime $\mathfrak r \subset S$ with $\mathfrak r \subset \mathfrak q$, see
Lemma \ref{lemma-Zariski-topology}.
By minimality of $\mathfrak p$ we see that
$\mathfrak p = \mathfrak r \cap R$.
\end{proof}

\begin{lemma}
\label{lemma-intersection-not-zero}
Let $A \subset B$ be an inclusion of domains inducing an algebraic extension
of fraction fields. If $J \subset B$ is a nonzero ideal, then $A \cap J$
is nonzero too. Thus the image of a proper closed subset of $\Spec(B)$
is not dense in $\Spec(A)$.
\end{lemma}

\begin{proof}
Let $x \in J$ be a nonzero element. Since $x$ is algebraic over the fraction
field of $A$, there exists a $d \geq 1$ and $a_0, \ldots, a_d \in A$
with $a_0, a_d \not = 0$ such that
$a_d x^d + a_{d - 1} x^{d - 1} + \ldots + a_0 = 0$ in $B$.
Then $a_0 \in A \cap J$.
\end{proof}




